De schapen in de wei
keken ineens opzij.

Kwam de slachter eraan?
Waarheen moesten ze gaan?

Werden ze weer verdeeld?
Naar klasse ingedeeld?

Deze slachter keek slechts
rustig van links naar rechts

zonder in te grijpen.
Puur om te begrijpen

wat de schapen echt zijn.
Ze vormden een gordijn,

die hij wilde doorzien.
Zijn kudde overzien.

In zijn meditatie
meldde een profetie

“wordt een wei voor dwazen,
die jouw boodschap grazen.”

“Ga je me verkrachten?”,
stond Svenska te wachten,

op een vertrouwde plek,
geketend aan het hek.

De profeet schudde nee.
Wat moest ze daar nu mee?

“Hou je van een vulva?”
De profeet knikte ja.
“Ik ben een slappe griet.
Waarom neem je me niet?”

“Ik geef me niet zomaar
aan iedere zondaar.

Zou je dat graag willen?”
Svenska wilde gillen

dat hoorde bij het spel.
Maar hij bleek een rebel,

die ze wilde kennen.
Misschien wel verwennen.

“Staat mijn rauwe argwaan
volgens jouw gelijk aan

idiotologie
of ideologie?”

Het werd hem zonneklaar
toen Svenska hem zowaar

het eerste antwoord gaf.
Aldus groeide zijn staf.

Hij die wilde sproeien
liet haar in de boeien.

Al bood ze geen verzet,
toch voorspelde hij het.

Zelfmoord

Huisvoogd of vader werd soms Gé Klungel of geklungel genoemd terwijl hij denkt dat ie wordt vergeleken met een andere werkelijke persoon

Emma schreef in een brief over zelfmoord nadat ze had ontdekt wie haar echte vader was. Dat schreef ze Emmy die door vragen te stellen en na te denken een verband had gelegd. Daarna heeft ze nooit meer op een brief geantwoord.

De wens om dood te willen
schoof door zichzelf te killen
van radicaal levendig
naar koelbloedig bestendig.

Stel je voor dat je elke laag van de werkelijkheid zou kunnen afpellen – de melkwegen, de sterren, de atomen, de subatomaire deeltjes – tot er niets meer overblijft dan het pure zijn van het bestaan zelf. Wat zou je dan aantreffen?

In de dertiende eeuw deed de filosoof en theoloog Thomas van Aquino een buitengewone uitspraak. God is niet eenvoudigweg het grootste wezen in het universum, noch één ding onder velen; God ís het zijn zelf, het pure bestaan zonder welk niets zou kunnen zijn. In het Latijn noemde hij dit “ipsum esse subsistens”: het subsisterende zijn, het bestaan zelf. Dit was geen poëtische uitdrukking of theologische metafoor. Voor Van Aquino was dit de meest accurate manier om de ultieme werkelijkheid te beschrijven die alles, overal, op dit moment in stand houdt. En hier komt het deel dat alles verandert: het zijn zelf is geen koude, onpersoonlijke kracht, maar bewust, oneindig gewaar en volkomen levend.

In deze video wordt de visie van Van Aquino in detail verkend, vanaf de wortels in de antieke filosofie tot aan de opmerkelijke overeenkomsten met de moderne wetenschap. We zullen bovendien ontdekken wat het betekent om te leven in een werkelijkheid waarin het bestaan zelf bewust is en goddelijk.

Om Van Aquino's inzicht echt te begrijpen, moeten we ons verplaatsen in zijn wereld: een Europa uit de dertiende eeuw, bruisend van de herontdekking van antieke kennis. Thomas van Aquino werd in 1225 geboren in Roccasecca, Italië, in een tijd dat de werken van Aristoteles via vertalingen van Arabische geleerden als Avicenna en Averroës opnieuw hun weg vonden naar het westerse denken. Dit was revolutionair; Aristoteles’ filosofie bood een nauwkeurige manier van denken over natuur, causaliteit en het bestaan zelf.

Als dominicaanse monnik was Van Aquino toegewijd aan het christelijk geloof, maar hij geloofde dat geloof en rede geen vijanden waren: waarheid kon waarheid niet tegenspreken. Door de metafysica van Aristoteles te verweven met de christelijke theologie ontwikkelde Van Aquino een visie op God die heel anders is dan waar de meeste mensen aan denken. God is niet een kosmische ambachtsman die aan de wereld sleutelt van buitenaf, maar de reden dát er een wereld is. Door Aristoteles’ onderscheiding tussen essentie – wat iets is – en bestaan – dát iets is – concludeerde Van Aquino dat alles wat wij zien zijn bestaan aan iets anders ontleent, en dat alleen God het zijn zelf is: onverklaard, noodzakelijk, oneindig.

Voor hem was dit niet alleen filosofie, maar ook een spirituele visie op de werkelijkheid als een onafgebroken acte van goddelijke aanwezigheid. Eén van Van Aquino’s belangrijkste correcties op hoe er gewoonlijk wordt gedacht, is dat God geen zijnde is – zelfs niet het machtigste zijnde dat bestaat – maar de daad van het zijn zelf. Een zijnde, hoe immens ook, blijft een eindig iets, met grenzen, dat samen bestaat met andere zijn-dingen in het grote inventaris van de werkelijkheid. Maar het zijn zelf is geen onderdeel van deze lijst – het is de reden dat die lijst überhaupt bestaat. Denk hierbij aan het verschil tussen afzonderlijke golven en de oceaan: de golven rijzen op en verdwijnen weer, maar het water, de oceaan zelf, blijft de voortdurende bron van hun bestaan. Van Aquino ging nog verder: zonder God zou er geen oceaan zijn, geen golven, geen water, geen zijn.

Daarom beschreef hij God als het subsisterende bestaan: puur zijn dat uit zichzelf bestaat, zonder oorsprong, zonder grenzen. Ander dan het blinde, mechanische ‘zijn’ van de abstracte filosofie is de God van Van Aquino pure act, “actus purus”: volledig gewaar, volledig kennend, volmaakt liefhebbend. Dit is het oneindige “Ik ben” die tot Mozes sprak vanuit de brandende braamstruik; niet gebonden aan ruimte of tijd, maar degene in wie alle tijd en ruimte überhaupt bestaan.

Als God het zijn zelf is, en het zijn zelf bewust is, dan is ons eigen bewustzijn geen geïsoleerd fenomeen gevangen in het brein, maar een vonk van die oneindige geest die het universum draagt. Van Aquino zag alle eindige bewustzijnsvormen als deelnames aan het goddelijke bewustzijn – zoals zonnestralen uitingen van de zon zijn. Jouw vermogen om “ik ben” te zeggen is niet gescheiden van het grote “Ik ben” dat het hart van de werkelijkheid vormt.

Dit idee heeft een diepe weerklank met mystieke tradities wereldwijd. In advaita vedanta is Brahman het pure bewustzijn waarin alle vormen verschijnen; in het soefisme is God het licht van de hemelen en de aarde; en in bepaalde moderne filosofische stromingen, zoals bij David Bentley Hart, wordt betoogd dat bewustzijn niet kan opstaan uit het absolute niets, maar verankerd moet zijn in de diepste laag van de werkelijkheid.

Voor Van Aquino zijn bestaan en bewustzijn onafscheidelijk; bestaan is alvast deelnemen aan weten en gekend worden. Dat wil zeggen, elk flikkeren van bewustzijn, elk moment waarop je beseft dat je bestaat, is al een subtiele ontmoeting met het goddelijke.

Van Aquino is waarschijnlijk het beroemdst om zijn vijf wegen – filosofische argumenten voor het bestaan van God. De meest relevante in deze context is het argument van de contingentie. Dit begint met een simpele observatie: alles wat we ervaren bestaat, maar had er ook niet kunnen zijn. Een boom groeit, maar had ook als zaad kunnen verrotten; een ster brandt, maar had nooit kunnen ontbranden. Deze contingente dingen zijn afhankelijk van iets anders voor hun bestaan. Als je die keten van afhankelijkheid terugvoert, kan dat niet oneindig doorgaan. Er moet iets zijn dat noodzakelijk bestaat – zonder afhankelijk te zijn van iets anders.

Cruciaal voor Van Aquino is dat deze noodzakelijke werkelijkheid niet zomaar een ander opperwezen aan de top van het universum is, maar de daad van het zijn zelf: onverklaard, eeuwig, de bron waaruit alle werkelijkheid stroomt. Moderne kosmologie raakt aan deze vraag als zij zich afvraagt: waarom is er iets en niet niets? Volgens Van Aquino komt het antwoord hierop neer dat op elk moment de werkelijkheid in het bestaan wordt geschonken door degene wiens essentie is te zijn.

Hoewel Aquino eeuwen vóór de kwantumfysica leefde, vertoont zijn visie opmerkelijke parallellen met de moderne wetenschap. In de kwantummechanica zijn deeltjes geen solide objecten met vaststaand bestaan; ze ontstaan door interactie, in een dans van waarschijnlijkheden die in realiteit samenkomen. De werkelijkheid is geen stilstaand blok, maar een voortdurend proces van worden, een continu scheppend gebaar dat aansluit bij Van Aquino’s idee van bestaan als doorlopende gave. In de neurowetenschap en de filosofie van de geest worstelen velen om bewustzijn te herleiden tot een bijverschijnsel van materie. Sommigen wenden zich tot het panpsychisme, waarin wordt gesteld dat bewustzijn een fundamenteel kenmerk van het universum is. Dit lijkt sterk op de gedachte bij Van Aquino dat bewustzijn in het weefsel van het zijn is verweven, omdat de bron van het zijn zelf bewust is.

Zelfs in de kosmologie, wanneer geleerden als Stephen Hawking zich afvragen wat leven blaast in de vergelijkingen van de fysica, raken ze aan hetzelfde mysterie dat Van Aquino beschreef: de eeuwige daad die werkelijkheid – hier en nu – in stand houdt.

Voor Van Aquino was dit nooit slechts een intellectuele puzzel, maar de basis voor een manier van leven. Als God het zijn zelf is, dan bestaat elke ademhaling, elke hartslag, elke gedachte alleen maar omdat het goddelijke ze in stand houdt. Dit betekent dat je niet door kosmische afstand van God gescheiden bent. Je bent nu, op dit moment, ondergedompeld in Gods werkelijkheid. Gebed, meditatie en contemplatie zijn daarom niet bedoeld om een verre God te bereiken, maar om je bewust te worden van de aanwezigheid die dichter bij je is dan je eigen hartslag.

En als elk schepsel deelheeft aan deze ene act van zijn, dan is elk mens, elk levend wezen, elk stukje van het universum een unieke uitdrukking van diezelfde oneindige bron. De ethische gevolgen hiervan zijn diepgaand: een ander kwaad doen is schade toebrengen aan het gedeelde weefsel van de werkelijkheid zelf; een ander liefhebben is de goddelijke aanwezigheid eren die jullie beiden in stand houdt.

De woorden van de apostel Paulus: “In Hem leven wij, bewegen wij en zijn wij” vatten de kern van Van Aquino’s visie samen. God is niet ergens anders, die af en toe in het universum ingrijpt. God is de reden dat het universum überhaupt bestaat. De plek waar je nu bent, het bewustzijn dat je op dit moment ervaart, het feit dat er überhaupt iets is in plaats van niets – dat alles is een uitdrukking van dezelfde ongeschapen act van zijn. Dit realiseren, betekent verschuiven van een beeld van God als een verre, geloofwaardige entiteit naar de ervaring van God als de levende werkelijkheid waarin je al verblijft.

Dit is het mysterie waartoe Van Aquino ons uitnodigt: het bestaan zelf als het bewuste, goddelijke Zijn. Leven in dit besef is leven in verwondering, dankbaarheid en ontzag – niet voor een God die concurreert met de schepping, maar voor een God die het leven van de schepping zélf is.

op dit moment, terwijl je deze video bekijkt, ervaar je iets dat niet bestaat. althans niet zoals wij denken dat het bestaat. het levendige rood van een appel, het diepe blauw van de lucht buiten je raam, de warme gouden gloed van je lamp: geen van deze kleuren bestaat in de wereld om je heen. het zijn biologische uitvindingen, vertalingen die volledig binnen de architectuur van je brein worden gecreëerd. in de objectieve realiteit die de fysica beschrijft, is er geen rood, er is geen blauw, er is helemaal geen kleur. er zijn alleen golflengten van elektromagnetische straling die door de ruimte stuiteren, onzichtbare, kleurloze trillingen die je zenuwstelsel transformeert in de rijke, verzadigde wereld die jij thuis noemt. en hier is wat je vannacht wakker zou moeten houden: als iets zo fundamenteels als kleur een complete fabricatie van je geest is, wat nog meer? hoe zit het met de solide objecten om je heen? hoe zit het met de ruimte ertussen? hoe zit het met de zeer eigen sensatie van jij-zijn, hier zitten in dit moment, terwijl deze woorden voor je ogen verschijnen? kijken we naar de wereld door een helder raam en observeren we de realiteit zoals die werkelijk is, of leven we binnen een minutieus gecomponeerd verhaal, moment na moment geschreven door de meest geavanceerde verhalenverteller die je ooit bent tegengekomen, je eigen geest?
dit is geen filosofie. dit is geen speculatie. dit is de verbluffende conclusie die opkomt uit de voorhoede van neurowetenschap en bewustzijnsonderzoek. wetenschappers ontdekken dat je brein de wereld om je heen niet passief ontvangt. het construeert die actief en schildert de realiteit van binnenuit met een palet van herinneringen, voorspellingen en verborgen aannames die je nooit bewust hebt onderzocht. de implicaties zijn verbijsterend. als je waarneming geen getrouwe opname is maar een actieve creatie, dan kan de fundamentele aard van je ervaring, alles wat je ooit hebt gezien, gehoord, gevoeld of geloofd over de wereld, iets totaal anders zijn dan je denkt. we gaan de verborgen architectuur van je eigen bewustzijn verkennen, de geheime machinerie die jouw realiteit bouwt, moment na moment.
deze reis is geworteld in toonaangevende wetenschap en filosofie, niet bedoeld om makkelijke antwoorden te geven, maar om je een nieuwe lens te bieden om het meest intieme mysterie van allemaal te bekijken: je eigen ervaring. wat we samen zullen ontdekken kan veranderen hoe je alles ziet. het kan veranderen hoe je jezelf ziet. als je gelooft dat de grootste mysteries de mysteries zijn die we in ons meedragen, overweeg dan te abonneren.
laten we beginnen met het vinden van de eerste barst in de spiegel van de perceptie, de plek waar de mooie misleiding van je brein zichzelf begint te verraden. op dit moment zit er een gat in je zicht. een blinde plek waarin je helemaal niets ziet. en je brein verbergt die al je hele leven voor je. sluit je rechteroog, houd je vinger op armlengte voor je en beweeg hem langzaam naar links terwijl je recht vooruit blijft kijken. op een gegeven moment zal je vinger gewoon verdwijnen. niet omdat hij te snel bewoog, niet omdat het licht veranderde, maar omdat je je blinde vlek vond, een gebied waar je oogzenuw op je netvlies aansluit, waardoor er een gat ontstaat in je gezichtsveld ter grootte van negen volle manen. en het verontrustende deel is: je merkt dit gapende gat in je zicht nooit op. je brein laat het niet leeg. het vult het niet met ruis of duisternis. het fabriceert actief wat er zou moeten zijn. als je naar een blauwe muur kijkt, schildert het blauw in de leegte. als er een patroon is, zet het dat patroon naadloos over de kloof voort. dit heet perceptuele invulling. en het is geen gok, geen benadering. je brein produceert daar letterlijk realiteit, creëert visuele informatie die nooit door je ogen is vastgelegd, en overtuigt je vervolgens zo volledig dat je nooit aan de misleiding twijfelt.
maar de montage gaat dieper. elke paar seconden schieten je ogen in snelle bewegingen rond, zogenaamde saccades. tijdens deze bliksemsnelle sprongen doet je brein iets buitengewoons: het verwijdert actief de binnenstromende visuele informatie. het maakt je doelbewust blind, het wist de duizelingwekkende waas die je zintuigen anders zou bestoken terwijl je ogen van punt naar punt springen. je bent functioneel blind tientallen keren per minuut, maar je ervaart nooit blindheid. je brein naait een naadloze, stabiele wereld aan elkaar uit deze losstaande momentopnames en verbergt zijn eigen machinerie zo compleet dat jij continue, ononderbroken visie ervaart.
denk na over wat dit betekent. je brein registreert niet simpelweg passief wat er buiten is; het monteert de feed in realtime. het neemt leidinggevende beslissingen over wat jij mag ervaren. het verwijdert informatie, fabriceert informatie en strijkt de gaten zo vakkundig glad dat jij je totaal niet bewust bent van het proces. filosoof daniel dennett realiseerde zich de diepe implicaties van deze simpele fysiologische feiten. als je brein je visuele ervaring moment na moment actief construeert, gaten invult en redactionele keuzes maakt, dan is wat jij zien noemt eigenlijk geen zien. het lijkt meer op gecontroleerd dromen terwijl je wakker bent. maar als je brein al je meest fundamentele zintuig redigeert, het zintuig dat je boven alle andere vertrouwt, wat redigeert het nog meer? welke andere aspecten van je realiteit worden actief geconstrueerd, gefilterd en aan jou gepresenteerd alsof ze gewoon daarbuiten lagen te wachten om ontdekt te worden? de blinde vlek is niet zomaar een anatomische curiositeit. het is een raam op de verborgen waarheid over perceptie zelf. jouw ervaring is geen passieve opname van de wereld. het is een scheppingsdaad, een realtime bouwproject binnen je schedel. en dit is nog maar het begin.
alles wat je is verteld over hoe je brein werkt, staat op z’n kop. eeuwenlang dachten we over het brein als een geavanceerde camera die signalen uit de wereld passief ontvangt en die samenstelt tot onze ervaring. licht raakt je ogen, geluidsgolven betreden je oren, en ergens in je schedel worden die signalen verwerkt tot de film van de realiteit waar je nu naar kijkt. maar dit is volledig verkeerd. je brein wacht niet tot de wereld het vertelt wat er gebeurt. het genereert voortdurend voorspellingen, gedetailleerde hypotheses van moment tot moment over wat er daarbuiten zou moeten zijn. nog voordat licht je netvlies raakt, heeft je brein al besloten wat je waarschijnlijk ziet. voordat geluid je oren bereikt, voorspelt het al wat je waarschijnlijk zult horen. dit is het revolutionaire kader dat neurowetenschappers predictive processing noemen, voorspellende verwerking, en het verandert alles wat we dachten te weten over perceptie.
stel je voor dat je in de schemering door een bos loopt. de zintuiglijke informatie die je brein bereikt is schaars en ruisig: schaduwen, geritsel, flarden van beweging in je perifere zicht. een camera zou moeite hebben om deze chaotische data te duiden. maar jouw brein heeft geen moeite, omdat het niet bij nul begint. het draait al een gedetailleerde simulatie van wat een bos in de schemering zou moeten bevatten: bomen die in de wind wiegen, kleine dieren die door struikgewas bewegen, het spel van wegstervend licht door bladeren. je zintuigen bouwen deze ervaring niet; ze corrigeren en verfijnen slechts een voorspelling die al draaide. daarom herken je de stem van je vriend in een druk restaurant, zelfs als je hem of haar nauwelijks kunt horen. je brein voorspelde dat die zo zou klinken en gebruikte net genoeg auditieve informatie om de gok te bevestigen. daarom kun je je slaapkamer in volledige duisternis navigeren: je brein draait een gedetailleerd ruimtelijk model van waar alles zou moeten staan en gebruikt slechts kleine snippers sensorische feedback om zijn voorspellingen te verifiëren.
neurowetenschapper karl friston, een van de architecten van deze theorie, beschrijft je brein als in de kern een voorspelmachine. maar misschien komt de meest elegante duiding hiervan van neurowetenschapper anil seth, die je hele perceptuele ervaring een gecontroleerde hallucinatie noemt. denk eens na over die uitdrukking. een hallucinatie is wanneer je brein ervaringen genereert die niet overeenkomen met de externe realiteit. maar wat seth suggereert is veel radicaler: alle perceptie is hallucinatie. het enige verschil is dat normale perceptie een hallucinatie is die gestuurd en gecorrigeerd wordt door zintuiglijke input. je ervaart niet de wereld, je ervaart de best mogelijke gok van je brein over de wereld, verfijnd door net genoeg zintuiglijke data om de hallucinatie in lijn met de realiteit te houden.
dit verklaart alles wat we in het vorige deel zagen. je blinde vlek wordt ingevuld omdat je brein voorspelde wat er zou moeten zijn. je zicht blijft stabiel tijdens oogbewegingen omdat je brein een stabiele wereld voorspelde en de tegenstrijdige informatie van je bewegende ogen simpelweg negeerde. maar hier wordt het diepgaand. als je hele perceptuele ervaring een gecontroleerde hallucinatie is die door voorspellingen wordt gegenereerd, dan lost de grens tussen binnen en buiten op. wat jij beschouwt als de externe wereld, de kamer waarin je zit, de geluiden om je heen, zelfs je eigen lichaam, staat niet los van je bewustzijn. het ís je bewustzijn, actief geconstrueerd, moment na moment, door het meest geavanceerde realiteitsgeneratiesysteem in het bekende universum. je brein toont je de wereld niet; het creëert een wereld voor jou om in te bewonen. en zodra je dit begrijpt, begin je te beseffen dat de architect van je universum al die tijd in het volle zicht verborgen was. dat was jij. of beter: de immense onbewuste intelligentie die jij noemt jij, schildert de realiteit van binnenuit met penseelstreken waarvan je nooit wist dat ze bestonden.
maar dit roept een verontrustende vraag op. als jouw realiteit actief wordt geconstrueerd door voorspellingen, wat gebeurt er dan als die voorspellingen worden beïnvloed door krachten waarvan je je niet eens bewust bent? je hart klopt nu. je maag verteert. je longen zetten uit en krimpen. en je brein luistert naar elk fluistersignaal uit je lichaam en gebruikt deze interne signalen als ingrediënten in de realiteit die het voor jou construeert. deze verborgen zin heet interoceptie: het bewustzijn van je brein van de interne toestand van je lichaam. en wat de neurowetenschap ontdekt, zal fundamenteel veranderen hoe je je eigen ervaring begrijpt. die lichamelijke signalen zijn niet slechts achtergrondruis. ze zijn primaire data die je voorspellende brein gebruikt om de wereld te schilderen die jij waarneemt.
denk aan de laatste keer dat je echt angstig was. zag de wereld er anders uit? leek die neutrale opmerking van een vriend plotseling beladen met oordeel? leken schaduwen dreigender? dat waren geen truukjes van je verbeelding. dat was je brein dat de data van angst, verhoogde hartslag, gespannen spieren, stresshormonen gebruikte als informatie over de wereld zelf. wanneer je zenuwstelsel geactiveerd is, denkt je voorspellende brein niet: “ik voel me angstig in een neutrale wereld.” het concludeert: “dit moet een gevaarlijke wereld zijn die waakzaamheid vereist.” vervolgens construeert het je perceptuele ervaring dienovereenkomstig, door potentiële bedreigingen te benadrukken, ambigue sociale signalen te versterken en de omgeving letterlijk vijandiger te laten aanvoelen. dit is wat neurowetenschapper lisa feldman barrett affectief realisme noemt: het diepgaande idee dat je emoties geen reacties op de wereld zijn, maar actieve ingrediënten in hoe je brein de wereld construeert die jij ervaart.
als je verdrietig bent, lijken kleuren daadwerkelijk minder verzadigd. dat is niet metaforisch. je brein, dat de lage energie en chemische signalen van verdriet interpreteert, voorspelt een wereld met minder levendigheid en construeert je visuele ervaring overeenkomstig. als je uitgeput bent, lijkt die heuvel inderdaad steiler. je brein, dat uitgeputte energiereserves aanvoelt, past zijn ruimtelijke voorspellingen aan om de extra inspanning te weerspiegelen die je huidige lichamelijke staat vereist. als je verliefd bent, verschijnt het gezicht van die persoon voor jou echt mooier dan voor anderen. je brein, overstroomd met bindingshormonen en positieve voorspellingen, construeert een versie van hun gelaatstrekken die aansluit bij je interne emotionele landschap. dit is geen bias die binnensijpelt in een verder objectief systeem. zo werkt het systeem. je gevoelens zijn data over de wereld, net zo goed als lichtgolven of geluidsgolven data over de wereld zijn. je brein behandelt ze allemaal als even valide informatiebronnen bij het bouwen van je ervaring van moment tot moment. barretts onderzoek onthult iets buitengewoons: er is geen neutrale observatie. elke perceptie wordt gekleurd door de interne staat van de waarnemer. de hongerige ervaart voedsel anders dan de verzadigde. de eenzame ervaart sociale interacties anders dan de verbonden persoon. de zelfverzekerde ervaart uitdagingen anders dan de bange. dit verklaart waarom twee mensen exact dezelfde gebeurtenis kunnen zien en totaal verschillende ervaringen hebben van wat er is gebeurd. ze ervoeren niet dezelfde wereld gefilterd door verschillende perspectieven. ze ervoeren letterlijk verschillende werelden, geconstrueerd door breinen die met verschillende interne data opereerden. jouw realiteit wordt niet alleen beïnvloed door je emoties en lichamelijke toestanden; ze wordt er deels door gecreëerd. de grens tussen innerlijke en uiterlijke ervaring lost volledig op wanneer je inziet dat je gevoelens net zozeer deel uitmaken van het constructieproces als fotonen die je netvlies raken.
maar dit roept een diepe en onrustbarende vraag op. als je interne toestanden actief de realiteit vormgeven die je waarneemt, en als dit proces grotendeels onbewust is, hoeveel van wat jij de objectieve wereld noemt is dan in feite slechts een reflectie van patronen en voorspellingen die je nooit hebt onderzocht? wat als de structuur van de realiteit zoals jij die kent is gevormd door krachten die veel fundamenteler zijn dan je emoties van moment tot moment? als je brein realiteit construeert vanuit voorspellingen die door je interne toestanden worden beïnvloed, dringt zich één vraag op: waarom zou evolutie ons zo ontwerpen? waarom zouden we de wereld niet gewoon zien zoals die werkelijk is? het antwoord zal alles verbrijzelen wat je dacht te weten over de aard van de realiteit zelf.
cognitief wetenschapper donald hoffman onderzoekt deze vraag al decennia, en zijn conclusie is even elegant als angstaanjagend: evolutie heeft onze hersenen gevormd om de objectieve werkelijkheid voor ons te verbergen, omdat de waarheid zien ons zou doden. denk aan je computerscherm. je ziet kleurrijke pictogrammen, een prullenbak, mappen, documenten. je klikt op deze afbeeldingen om met je bestanden te werken. maar dit zie je niet: de miljoenen regels code, de spanningspatronen, de magnetische toestanden op je schijf die de echte werkelijkheid van je computer vormen. de desktopinterface verbergt die complexiteit doelbewust. als je elke keer ruwe code moest manipuleren om een document te openen, zou je nooit iets gedaan krijgen. de kleurrijke pictogrammen lijken in niets op de onderliggende realiteit, maar ze zijn oneindig veel bruikbaarder om te overleven in de digitale wereld. hoffman betoogt dat perceptie precies zo werkt. ruimte, tijd, objecten, kleuren: dit zijn onze evolutionaire desktopiconen. het zijn nuttige ficties die ons in staat stellen door de realiteit te navigeren zonder te worden overweldigd door haar ware complexiteit.
bedenk hoe de objectieve realiteit er volgens de fysica uitziet: kwantumvelden die fluctueren in waarschijnlijkheidswolken, deeltjes die gelijktijdig in meerdere toestanden bestaan, ruimtetijd die buigt en kromt, energie en materie als uitwisselbare aspecten van een onderliggende wiskundige structuur. een organisme dat de realiteit zo waarneemt, zou verlamd worden door oneindige complexiteit. maar een organisme dat vereenvoudigde iconen ziet, solide objecten die door lege ruimte bewegen in lineaire tijd, kan snel en resoluut handelen. het kan vluchten voor roofdieren, voedsel vinden en zich voortplanten. dit is geen speculatie; dit is de onvermijdelijke wiskundige uitkomst wanneer je de kosten van het verwerken van oneindige informatie afzet tegen de baten van handelen op vereenvoudigde, overlevingsrelevante aanwijzingen. wij zijn de nakomelingen van ontelbare generaties interfacegebruikers. elke voorouder die lang genoeg overleefde om zich voort te planten, was een voorouder die nuttige ficties zag in plaats van overweldigende waarheid. wij bestaan omdat onze afstamming zich specialiseerde in het construeren van behulpzame hallucinaties, niet in het waarnemen van accurate realiteit.
de ruimte waarin je zit bestaat niet. de tijd die ogenschijnlijk om je heen stroomt is een constructie. de solide objecten om je heen zijn even fictief als de desktopiconen op je scherm. zelfs je gevoel een afzonderlijk zelf te zijn, gelokaliseerd in een lichaam, is gewoon een ander interface-element: nuttig voor overleving, maar volledig losgezongen van de onderliggende aard van wat je werkelijk bent. dit is de meest diepgaande paradigmawisseling in de menselijke geschiedenis. wij zijn geen toeschouwers van een objectieve wereld. we zijn actieve deelnemers aan een collectieve hallucinatie, geëvolueerd over miljoenen jaren, zo overtuigend dat we haar verwarren met de realiteit zelf. je hebt de wereld nooit gezien. je hebt alleen ooit de versie van je geest ervan gezien. maar als dit waar is, als alles wat je ooit hebt ervaren slechts een interface is die een onkenbare realiteit verbergt, wat ben jij dan? en wat is de aard van het doek waarop je bewustzijn deze uitgebreide illusie schildert?
als alles wat je ooit hebt ervaren slechts een interface is die een onkenbare realiteit verbergt, kun je een diep gevoel van ontkoppeling voelen, verloren, afgesneden van de waarheid van het bestaan zelf. maar wat als dit gevoel slechts een andere voorspelling is die je brein maakt? en wat als die voorspelling fout is? jij bent geen passief slachtoffer van een uitgekiende misleiding. je bent geen gevangene achter de sluier van de perceptie, voor altijd gescheiden van de echte wereld. je bent iets veel buitengewoners dan dat. je bent een bewust handelende agent die actief samenwerkt met het universum in de voortdurende creatie van de ervaring zelf. elke moment is je gewaarzijn verwikkeld in de meest intieme dans die denkbaar is met welke fundamentele realiteit er ook maar voorbij de interface ligt. je bewustzijn observeert niet simpelweg een reeds bestaande wereld en ontvangt geen vooraf bepaalde hallucinatie passief. het doet mee. het kiest. het creëert.
wanneer je je aandacht verlegt, verander je niet alleen wat je opvalt aan een vaste realiteit; je neemt deel aan de constructie, moment na moment, van wat voor jou echt wordt. wanneer je je interne toestand verandert via je adem, via beweging, via de kwaliteit van je gedachten, verander je niet slechts je stemming; je verandert het weefsel van de wereld die je bewoont. dit is geen filosofie. dit is geen wensdenken. dit is de diepgaande implicatie van alles wat we samen hebben verkend. als perceptie constructie is, en bewustzijn de constructeur, dan is gewaarzijn zelf de meest creatieve kracht die je ooit zult tegenkomen. de vraag is niet of je verbonden bent met de ultieme realiteit. de vraag is: wat is de aard van het mysterieuze doek waarop je bewustzijn schildert?
wat is het fundamentele substraat dat reageert op de penseelstreken van je aandacht, je intentie, je zijn? dit is de rand van het bekende, de plek waar neurowetenschap oplost in mysterie en mysterie opent naar verwondering. je geest creëert je realiteit. maar wat creëert je geest? je bewustzijn construeert je wereld. maar wat is bewustzijn zelf? je kwam hierheen in de veronderstelling dat je iets over perceptie zou ontdekken. in plaats daarvan stuitte je op de diepste vraag van allemaal: wat ben je werkelijk, wanneer alles waarvan je dacht dat je het was, een interface blijkt te zijn die iets onvoorstelbaar veel diepers verbergt? dit zijn de territoria die we samen verkennen, de vragen die leven op het snijvlak van wetenschap en mysterie, kennis en verwondering. beyond reality is nog maar net begonnen.
